% Part: first-order-logic
% Chapter: syntax-and-semantics
% Section: main-operator

\documentclass[../../../include/open-logic-section]{subfiles}

\begin{document}

\olfileid{fol}{syn}{mai}

\olsection{\printtoken{S}{main operator} van een formule}

\begin{explain}
Vaak willen we weten welk logisch teken als laatste is gebruikt om
!!a{formula}~$!A$ op te bouwen. Dat teken heet het \emph{buitenste
  logische teken} van~$!A$: het staat als het ware om de rest van
$!A$ heen. Het buitenste logische teken van $\lnot !A$ is bijvoorbeeld
$\lnot$. Bij $(!A \lor !B)$ is het $\lor$, enzovoort.
\end{explain}


\begin{defn}[!!^{main operator}]
\ollabel{def:main-op}
Zo bepalen we het \emph{!!{main operator}} van !!a{formula}~$!A$:
\begin{enumerate}
\item \indcase*{!A}{!A}{$\indfrm$ heeft geen !!{main operator}.}

\tagitem{prvNot}{\indcase{!A}{\lnot !B}{het !!{main operator} van $\indfrm$
  is~$\lnot$.}}{}

\tagitem{prvAnd}{\indcase{!A}{(!B \land !C)}{het !!{main operator} van
  $\indfrm$ is~$\land$.}}{}

\tagitem{prvOr}{\indcase{!A}{(!B \lor !C)}{het !!{main operator} van
  $\indfrm$ is~$\lor$.}}{}

\tagitem{prvIf}{\indcase{!A}{(!B \lif !C)}{het !!{main operator} van
  $\indfrm$ is~$\lif$.}}{}

\tagitem{prvIff}{\indcase{!A}{(!B \liff !C)}{het !!{main operator} van
  $\indfrm$ is~$\liff$.}}{}

\tagitem{prvAll}{\indcase{!A}{\lforall[x][!B]}{het !!{main operator}
  van $\indfrm$ is~$\lforall$.}}{}

\tagitem{prvEx}{\indcase{!A}{\lexists[x][!B]}{het !!{main operator} van
  $\indfrm$ is~$\lexists$.}}{}
\end{enumerate}
\end{defn}

We bedoelen telkens precies de aangegeven \emph{plaats waar het teken
voorkomt} in de formule. Kijk bijvoorbeeld naar de formule
$((!D \lif !E) \lif (!E \lif !D))$. Zij heeft de vorm $(!B \lif !C)$,
met $!B$ gelijk aan $(!D \lif !E)$ en $!C$ gelijk aan $(!E \lif !D)$.
Daarom is het tweede voorkomen van $\lif$ het !!{main operator}.

\begin{explain}
Dit is een \emph{recursieve} definitie: zij koppelt iedere niet-atomaire
!!{formula} aan de plaats waar haar !!{main operator} voorkomt. De
regels bouwen !!{formula}s stap voor stap op. Daarom valt iedere
!!{formula}~$!A$ onder een van de gevallen in \olref{def:main-op}.
Iedere niet-atomaire !!{formula}~$!A$ heeft dus !!a{main operator}.
Verder valt iedere !!{formula} onder maar één geval. Ook liggen de
kleinere !!{formula}s waaruit $!A$ is opgebouwd precies vast. Er is dus
maar één plaats die het !!{main operator} van~$!A$ kan zijn. Daarmee
hebben we inderdaad een functie bepaald.
\end{explain}

De naam van !!{formula}s hangt af van het symbool dat hun !!{main
operator} is. \olref{tab:main-op} geeft de namen.
\iftag{defNot,defOr,defAnd,defIf,defIff,defTrue,defFalse,defEx,defAll}
{ Gedefinieerde operatoren staan officieel niet in !!{formula}s. Het
zijn alleen afkortingen. Daarom kunnen zij officieel niet het buitenste
logische teken van een formule zijn. In bewijzen over alle !!{formula}s
hoeven we ze dus niet als aparte gevallen te behandelen.}

\begin{table}[!h]
\centering
\begin{tabular}{c | c | c}
!!^{main operator} & Soort !!{formula} & Voorbeeld\\
\hline
geen & !!{formula} zonder logische deelopbouw &
\iftag{prvFalse}{$\lfalse$,}{}
\iftag{prvTrue}{$\ltrue$,}{}
$\Atom{R}{t_1, \dots, t_n}$,
$\eq[t_1][t_2]$\\
$\lnot$ & ontkenning & $\lnot !A$ \\
$\land$ & en-verbinding & $(!A \land !B$) \\
$\lor$ & of-verbinding & $(!A \lor !B$) \\
$\lif$ & !!{conditional} & $(!A \lif !B$) \\
$\liff$ & !!{biconditional} & $(!A \liff !B)$ \\
$\lforall[][]$ & !!{formula} met ``voor alle''& $\lforall[x][!A]$ \\
$\lexists[][]$ & !!{formula} met ``er bestaat''& $\lexists[x][!A]$
\end{tabular}
\caption{Buitenste logische teken en namen van !!{formula}s}
\ollabel{tab:main-op}
\end{table}
\end{document}
